% =====================================================================
% Shreya Mishra — B.Sc. (Honors) Dissertation
% The Economic Impacts of Floods: Evidence from Eastern India
% Compile with: tectonic dissertation.tex
% =====================================================================
\documentclass[12pt,a4paper,oneside]{report}

\usepackage[utf8]{inputenc}
\usepackage[T1]{fontenc}
\usepackage{lmodern}
\usepackage{mathptmx}                  % Times-like body font
\usepackage[english]{babel}
\usepackage[a4paper,left=1.25in,right=1in,top=1in,bottom=1in]{geometry}
\usepackage{setspace}
\usepackage{graphicx}
\usepackage{booktabs}
\usepackage{array}
\usepackage{longtable}
\usepackage{tabularx}
\usepackage{caption}
\usepackage{amsmath,amssymb}
\usepackage{titlesec}
\usepackage{fancyhdr}
\usepackage{tocloft}
\usepackage{enumitem}
\usepackage{microtype}
\usepackage[hidelinks,bookmarks=true,bookmarksopen=true]{hyperref}

% --- Spacing ---
\onehalfspacing
\setlength{\parskip}{6pt}
\setlength{\parindent}{0pt}

% --- Caption style ---
\captionsetup{font={small,it},labelfont={bf},labelsep=colon,justification=centering,singlelinecheck=false}
\captionsetup[figure]{name=Figure}
\captionsetup[table]{name=Table}

% --- Chapter / section heading styles ---
\titleformat{\chapter}[display]
  {\normalfont\Large\bfseries\centering}
  {\MakeUppercase{\chaptertitlename} \thechapter}
  {12pt}
  {\Large\MakeUppercase}
\titlespacing*{\chapter}{0pt}{0pt}{20pt}

\titleformat{\section}{\normalfont\large\bfseries}{\thesection}{1em}{}
\titleformat{\subsection}{\normalfont\normalsize\bfseries\itshape}{\thesubsection}{1em}{}

% --- ToC formatting ---
\renewcommand{\cftchapfont}{\bfseries}
\renewcommand{\cftchappagefont}{\bfseries}
\setlength{\cftbeforechapskip}{6pt}

% --- Page styles ---
\fancypagestyle{plain}{%
  \fancyhf{}%
  \fancyfoot[C]{\thepage}%
  \renewcommand{\headrulewidth}{0pt}%
}
\pagestyle{plain}

% --- Custom commands for front matter "chapter*" pages ---
\newcommand{\frontchapter}[1]{%
  \clearpage
  \begin{center}\Large\bfseries\MakeUppercase{#1}\end{center}
  \vspace{1em}
  \addcontentsline{toc}{chapter}{\texorpdfstring{\MakeUppercase{#1}}{#1}}
}

% =====================================================================
\begin{document}

% =====================================================================
% TITLE PAGE
% =====================================================================
\begin{titlepage}
\begin{center}
\vspace*{0.5in}
{\Large\bfseries THE ECONOMIC IMPACTS OF FLOODS:\\[6pt]
EVIDENCE FROM EASTERN INDIA}\\[1.0in]

{\large A Thesis}\\[6pt]
{\normalsize Submitted in partial fulfillment of the requirements for the award of the Degree of}\\[6pt]
{\large\bfseries Bachelor of Science (Honors)}\\[6pt]
{\normalsize IN}\\[6pt]
{\large Quantitative Economics and Data Science}\\[0.6in]

{\normalsize BY}\\[6pt]
{\Large\bfseries SHREYA MISHRA}\\[4pt]
{\normalsize (IED/10032/23)}\\[1.5in]

{\normalsize\bfseries CENTRE FOR QUANTITATIVE ECONOMICS AND DATA SCIENCE}\\[4pt]
{\normalsize\bfseries BIRLA INSTITUTE OF TECHNOLOGY}\\[4pt]
{\normalsize\bfseries MESRA, RANCHI -- 835215}\\[4pt]
{\normalsize\bfseries 2026}
\end{center}
\end{titlepage}

% =====================================================================
% FRONT MATTER (roman page numbers)
% =====================================================================
\pagenumbering{roman}
\setcounter{page}{1}

% --- Approval of the Guide ---
\frontchapter{Approval of the Guide}

\noindent Recommended that the thesis entitled \textit{``The Economic Impacts of Floods: Evidence from Eastern India''} prepared by \textbf{Ms.\ Shreya Mishra} (Roll No.\ IED/10032/23) under my supervision and guidance be accepted as fulfilling this part of the requirements for the award of the Degree of Bachelor of Science (Honors) in Quantitative Economics and Data Science. To the best of my knowledge, the contents of this thesis did not form a basis for the award of any previous Degree to anybody else.

\vspace{0.6in}
\noindent Date: \rule{2in}{0.4pt}

\vspace{1.5in}
\begin{flushright}
\textbf{Dr.\ Saheli Bose}\\
Assistant Professor\\
Centre for Quantitative Economics and Data Science\\
Birla Institute of Technology\\
Mesra, Ranchi
\end{flushright}

% --- Declaration Certificate ---
\frontchapter{Declaration Certificate}

\noindent I certify that
\begin{enumerate}[label=(\alph*),itemsep=2pt]
\item The work contained in the thesis is original and has been done by myself under the general supervision of my supervisor.
\item The work has not been submitted to any other Institute for any other degree or diploma.
\item I have followed the guidelines provided by the Institute in writing the thesis.
\item I have conformed to the norms and guidelines given in the Ethical Code of Conduct of the Institute.
\item Whenever I have used materials (data, theoretical analysis, and text) from other sources, I have given due credit to them by citing them in the text of the thesis and giving their details in the references.
\item Whenever I have quoted written materials from other sources, I have put them under quotation marks and given due credit to the sources by citing them and giving required details in the references.
\end{enumerate}

\vspace{1.2in}
\begin{flushright}
\textbf{Shreya Mishra}\\
(IED/10032/23)
\end{flushright}

% --- Certificate of Approval ---
\frontchapter{Certificate of Approval}

\noindent This is to certify that the work embodied in this thesis entitled \textit{``Economic Impacts of Floods: Evidence from Eastern India''} is carried out by \textbf{Ms.\ Shreya Mishra} (Roll No.\ IED/10032/23) has been approved for the award of the Degree of Bachelor of Science (Honors) in Quantitative Economics and Data Science by the Birla Institute of Technology, Mesra, Ranchi.

\vspace{0.6in}
\noindent Date: \rule{2in}{0.4pt}\\[6pt]
Place: \rule{2in}{0.4pt}

\vspace{1.0in}
\noindent\begin{tabular}{p{0.45\textwidth}p{0.45\textwidth}}
\centering Internal Examiner & \centering External Examiner \tabularnewline
\end{tabular}

\vspace{0.8in}
\begin{center}
\textbf{Head of Department}
\end{center}

% --- Acknowledgement ---
\frontchapter{Acknowledgement}

I would like to express my profound gratitude to my project guide, \textbf{Dr.\ Saheli Bose}, for her invaluable guidance, constant encouragement, and constructive criticism throughout the course of this research. Her insights into empirical methodology and her patience with my iterative drafts have shaped this thesis at every stage.

I convey my sincere gratitude to the Head, Centre for Quantitative Economics and Data Science, Birla Institute of Technology, Mesra, for providing me with the opportunity to undertake this study and for the academic environment that made it possible.

I would also like to thank all the faculty members of the department who have directly or indirectly contributed to my understanding of econometrics, panel data, and applied economic research. The clarity I now have on fixed-effects identification owes a great deal to their teaching.

Finally, I must express my very profound gratitude to my parents for providing me with unfailing support and continuous encouragement throughout my years of study and through the process of researching and writing this thesis. This accomplishment would not have been possible without them.

Thank you.

\vspace{1.2in}
\noindent
\begin{tabular}{@{}p{0.45\textwidth}p{0.45\textwidth}@{}}
DATE: \rule{1.5in}{0.4pt} & \begin{flushright}\textbf{Shreya Mishra}\\(IED/10032/23)\end{flushright} \tabularnewline
\end{tabular}

% --- Abstract ---
\frontchapter{Abstract}

Floods are India's most recurrent natural disaster, and Eastern India --- comprising Bihar, Jharkhand, Odisha, and West Bengal --- bears a disproportionate share of the national flood burden. While the human and physical costs of these events are well documented, the local-level economic consequences remain difficult to quantify because conventional district output data are released with long lags and rarely capture short-run shocks. This study addresses that gap by using satellite-derived proxies to estimate the short-run economic impact of flood exposure on Assembly Constituencies (ACs) in Eastern India over 2014--2019.

The dependent variables are the year-on-year change in log nighttime lights (mean and median radiance from VIIRS), used as a proxy for local economic activity, and the Normalized Difference Built-up Index (NDBI, mean and median, from Landsat-8), used as a proxy for the stock of built-up infrastructure. Flood exposure is measured by the seasonal-water area share of each AC, constructed from the JRC Global Surface Water dataset. The empirical specification regresses the contemporaneous outcome on flood exposure and lagged values of NDVI, NDBI, and (in the NDBI regressions) nighttime lights, with Assembly Constituency and year fixed effects, and standard errors clustered at the district level. Each model is estimated both pooled across the four states and separately for Bihar, Jharkhand, Odisha, and West Bengal. The estimation sample contains 765 ACs across 108 districts and 3{,}890 AC-year observations.

The pooled regression yields a flood coefficient that is statistically indistinguishable from zero, but state-by-state estimation reveals sharp heterogeneity. In West Bengal, flood exposure significantly reduces nighttime-lights growth, with coefficients ranging from $-0.32$ (median, $p<0.05$) to $-0.39$ (mean, $p<0.05$). In Bihar, flood exposure significantly reduces the built-up index ($\beta_1 \approx -0.02$, $p<0.01$) and shows a marginal negative effect on nighttime lights. Jharkhand and Odisha display no consistently significant effects. The lagged dependent variable in the NDBI specification is subject to a small Nickell bias, but the contemporaneous flood coefficient --- the parameter of interest --- is largely unaffected. These findings indicate that flood impacts in Eastern India operate through state-specific channels --- economic activity in West Bengal, structural built-up in Bihar --- and suggest that uniform flood-mitigation policy is unlikely to be effective.

\vspace{0.2in}
\noindent\textbf{Keywords:} floods, nighttime lights, NDBI, satellite remote sensing, district panel, Eastern India, fixed-effects estimation.

% --- Table of Contents ---
\clearpage
\renewcommand{\contentsname}{Contents}
\tableofcontents

% --- List of Figures ---
\clearpage
\addcontentsline{toc}{chapter}{List of Figures}
\listoffigures

% --- List of Tables ---
\clearpage
\addcontentsline{toc}{chapter}{List of Tables}
\listoftables

% =====================================================================
% MAIN MATTER (arabic page numbers)
% =====================================================================
\clearpage
\pagenumbering{arabic}
\setcounter{page}{1}

% =====================================================================
\chapter{Introduction}
% =====================================================================

\section{Background}
India experiences a wide range of natural disasters every year, but few are as recurrent or as economically consequential as floods. According to the National Disaster Management Authority (NDMA, Government of India), floods account for the largest share of disaster-related damage in the country, both in terms of physical destruction and population affected. The National Flood Commission estimates that approximately 40 million hectares of land in India are flood-prone, of which around 7.2 million hectares are inundated each year. The Central Water Commission (CWC) places the average annual loss to crops, housing, and public utilities at roughly Rs.~1{,}805 crore, although these figures are widely understood to underestimate the true cost because they exclude indirect losses such as lost workdays, supply-chain disruption, and long-run capital depreciation.

The frequency and intensity of flood events have risen over the past two decades, driven by changing monsoon patterns, glacial melt in the Himalayan headwaters, embankment failures, and unplanned urbanisation in low-lying areas. As climate variability accelerates, the gap between the speed of flood events and the speed at which administrative damage assessments are released has widened. By the time conventional district-level output statistics become available, the policy window for response has often closed.

The present study takes this measurement gap as its starting point. Rather than relying on lagged administrative records, we use satellite-derived proxies --- nighttime lights, the Normalized Difference Built-up Index (NDBI), and the Normalized Difference Vegetation Index (NDVI) --- to construct an annual panel of economic outcomes at the Assembly Constituency level for the four most flood-affected states of Eastern India.

\section{Why Floods?}
Among all natural disasters that affect India, floods are the most frequent, the most spatially extensive, and the most damaging in absolute economic terms. Their importance for empirical economic analysis rests on three features. First, they are recurrent rather than exceptional, which provides repeated within-unit variation that can be exploited for identification. Second, they affect both physical capital (housing, roads, structures) and natural capital (cropland, soil), which means their impact propagates through several channels of the local economy. Third, they are observable from space at high temporal frequency, which solves the data-availability problem that has historically constrained sub-national flood research in India.

Despite this, empirical work on the economic consequences of flooding in India remains thin. Most existing studies are descriptive, focus on a single river basin, or are conducted at the state or district level --- too coarse to capture the heterogeneity that policy makers face on the ground. The granularity adopted in this study (Assembly Constituency $\times$ year) is finer than what most published Indian flood-economics research currently offers.

\section{Why Eastern India?}
Eastern India --- Bihar, Jharkhand, Odisha, and West Bengal --- is the most flood-affected region of the country, but each state experiences flooding through a distinct hydrological and economic profile.

\textbf{Bihar} lies at the centre of the Indo-Gangetic plain and is drained by transboundary rivers --- the Kosi, Gandak, and Bagmati --- that originate in Nepal and discharge into the Ganga. These rivers carry heavy sediment loads from the Himalayas, raise riverbed levels behind embankments, and produce recurrent breach-induced flooding across northern districts. National risk assessments place Bihar's share of India's flood-prone area at approximately 17.2 per cent (NIDM/NRSC).

\textbf{West Bengal} sits at the lower end of the Ganga--Brahmaputra system, where deltaic geomorphology, tidal influence, and the cyclones of the Bay of Bengal combine to produce both fluvial and coastal flood risk. The Damodar, Hooghly, and Sundarbans estuarine system make districts such as Murshidabad, Howrah, and the Sundarbans highly exposed to compound flood events.

\textbf{Odisha} faces a similarly compound risk profile, but with an even stronger coastal component. The Mahanadi delta is exposed to fluvial flooding upstream and to storm-surge flooding from cyclones that make landfall regularly along the Bay of Bengal coast. Coastal districts such as Puri, Kendrapara, and Jagatsinghpur are repeatedly affected.

\textbf{Jharkhand} is topographically distinct --- drained by the Damodar, Subarnarekha, and Brahmani--Mahanadi tributaries --- but flooding remains a serious risk in low-lying areas of Sahebganj, Pakur, and parts of the Damodar basin, especially in years of intense monsoon discharge from upstream Jharkhand hills.

The economic vulnerability of these states compounds their physical exposure. All four rank below the national average on per-capita Net State Domestic Product (MoSPI), and a substantial share of their workforce remains in agriculture. A flood that destroys a season's standing crop is therefore not merely a localised shock; it transmits through farm incomes, food supply chains, and migration flows. For these reasons, Eastern India offers both the variation needed for identification and the policy relevance needed to make the findings useful.

\section{Motivation of Study}
The motivation for this study is both empirical and policy-oriented. Empirically, floods generate sharp, plausibly exogenous shocks to local economic activity that can be observed at fine spatial and temporal resolution using remote-sensing data. This makes them a natural setting for fixed-effects panel estimation: the unit of observation is fixed (the Assembly Constituency), the shock varies over time within each unit, and confounders such as long-run development and common annual macro-shocks can be absorbed by AC and year fixed effects respectively.

The policy motivation is more direct. Flood mitigation in Eastern India currently relies on engineering interventions --- embankments, barrages, and drainage works --- designed and budgeted at the state level. Whether these uniform approaches are appropriate depends on whether floods damage Bihar, Jharkhand, Odisha, and West Bengal through the same channels. If they do, a centralised policy may suffice; if they do not, state-tailored responses are needed. Distinguishing these two scenarios empirically is the core motivation of this study.

\section{Research Questions}
The study is organised around three research questions:
\begin{enumerate}[leftmargin=2em,itemsep=4pt]
\item \textbf{Do floods reduce economic activity in flood-affected Assembly Constituencies?} Economic activity is proxied by year-on-year growth in nighttime-light radiance. The hypothesis is that contemporaneous flood exposure produces a negative and statistically significant coefficient on the flood variable in a fixed-effects panel regression.
\item \textbf{Does the impact of floods on economic activity vary across the four states of Eastern India?} Estimating the model separately for Bihar, Jharkhand, Odisha, and West Bengal allows the flood coefficient to differ by state. A non-trivial pattern of state-level heterogeneity would imply that pooled regressions mask important variation.
\item \textbf{What is the impact of flood inundation on built-up infrastructure?} Built-up intensity is proxied by the Normalized Difference Built-up Index (NDBI). This question separates the flow effect of floods (loss of activity, captured by nighttime lights) from the stock effect (damage to physical structures, captured by NDBI).
\end{enumerate}

\section{Objectives}
Corresponding to the research questions, the objectives of this study are:
\begin{enumerate}[label=(\roman*),leftmargin=2.2em,itemsep=4pt]
\item To construct a balanced panel of Assembly Constituency $\times$ year observations for Bihar, Jharkhand, Odisha, and West Bengal over the period 2014--2019, combining flood-exposure data from the JRC Global Surface Water dataset with satellite-derived outcomes (nighttime lights, NDBI, NDVI).
\item To estimate the contemporaneous effect of flood exposure on the year-on-year change in nighttime-lights radiance, using OLS with Assembly Constituency and year fixed effects and standard errors clustered at the district level.
\item To estimate the corresponding effect on the level of NDBI, controlling for the lagged dependent variable and other lagged covariates.
\item To compare pooled and state-by-state estimates in order to characterise the cross-state heterogeneity of flood impacts across all four states of Eastern India.
\item To interpret the findings in light of the agronomic and infrastructural channels through which floods affect Eastern Indian economies.
\end{enumerate}

\section{Gaps Addressed}
This study contributes to four gaps in the existing literature on Indian flood economics:
\begin{enumerate}[label=(\roman*),leftmargin=2.2em,itemsep=4pt]
\item \textbf{Spatial granularity.} Most prior work measures flood impact at the state or district level; the AC-level panel used here is finer-grained and closer to the units that respond to flood events on the ground.
\item \textbf{Temporal frequency.} Survey-based output measures are released with long lags; satellite proxies (nighttime lights, NDBI, NDVI) provide annual observations available within months of the event.
\item \textbf{Flow versus stock.} Existing work seldom separates the activity (flow) effect of floods from the infrastructure (stock) effect; this study estimates both within a unified specification.
\item \textbf{State-level heterogeneity.} Almost no published Indian study compares the flood coefficient across Bihar, Jharkhand, Odisha, and West Bengal. This study makes that comparison its central object of analysis.
\end{enumerate}

\section{Thesis Organization}
The remainder of this thesis is organised as follows. Chapter~2 sets out a minimal theoretical framework linking flood exposure to observable economic outcomes through three channels --- productive capital, agricultural land, and economic disruption. Chapter~3 describes the data, panel structure, and summary statistics, including temporal trends. Chapter~4 develops the two estimating equations, the choice of fixed effects, and the estimation strategy. Chapter~5 presents the four main regression tables and supporting visuals (forest plot, coefficient heatmap, robustness check) and discusses the cross-state heterogeneity. Chapter~6 summarises the findings, draws policy implications, and outlines directions for future work. References and an appendix follow.

% =====================================================================
\chapter{Theoretical Framework}
% =====================================================================

\section{Introduction}
The research questions posed in Chapter~1 ask whether floods reduce local economic activity, whether the impact varies across the four states of Eastern India, and whether physical infrastructure is affected differently from economic flow. To answer these questions empirically, we first need a framework that specifies how a flood is expected to translate into observable economic outcomes. This chapter sets out that framework. Section~2.2 identifies three channels through which flood exposure plausibly affects the local economy. Section~2.3 explains why three satellite-derived indices --- nighttime lights, NDBI, and NDVI --- are the appropriate empirical counterparts to those channels.

\section{Channels Linking Floods to Local Economic Activity}
A flood event affects local economic activity through three principal channels, distinguished by the type of asset they impair and the time horizon over which the impairment is observable.

\textbf{(i) The productive-capital channel.} Standing infrastructure --- housing, commercial structures, transport networks, and electricity assets --- depreciates when inundated. Once water recedes, the residual stock is smaller than before. The reduction may be gradual (from accelerated wear) or abrupt (from collapse), but in either case the post-flood capital stock supports a lower level of nighttime activity. This channel implies that flood exposure should be associated with a fall in observed economic-activity proxies that depend on installed capital, particularly nighttime-light radiance.

\textbf{(ii) The agricultural-land channel.} Inundation directly damages standing crops, salinises soil, and disrupts the planting calendar for the subsequent season. In Eastern India, where a substantial share of the workforce remains in agriculture, this channel transmits the flood shock from farm income through to local consumption, retail demand, and informal economic activity. The agronomic impact has both a short-run negative component (lost output of the affected season) and, in some settings, a positive long-run component arising from alluvial silt deposition and groundwater recharge --- a pattern documented for Indian river basins by Kale (2003).

\textbf{(iii) The economic-disruption channel.} Floods interrupt mobility, supply chains, and labour markets even where physical capital is not destroyed. Roads become impassable, markets close temporarily, and workers are displaced. These effects are transitory but contemporaneous with the flood event, which is why the empirical specification places flood exposure at time $t$ on the right-hand side rather than as a lag.

The three channels are not mutually exclusive. A single flood typically engages all three, but their relative magnitude varies with the structural composition of the local economy. In a predominantly agrarian district, the agricultural channel dominates; in a built-up urban area, the productive-capital and disruption channels matter more. This compositional dependence is one reason why flood impacts may differ across the four states of Eastern India even when the physical exposure (the share of land inundated) is comparable.

\section{Why Satellite Proxies?}
Direct measurement of the three channels requires data on output, capital stock, agricultural production, and trade flows at sub-district frequency. No such dataset exists for India at the Assembly Constituency level. We therefore proxy the unobservable channels with three satellite-derived indices, each chosen because it captures a specific dimension of the local economy.

\textbf{Nighttime lights (NL).} Annual VIIRS Day--Night Band radiance provides a well-validated proxy for local economic activity. A drop in NL between consecutive years indicates either reduced electricity consumption (consistent with the disruption channel) or a smaller stock of lit structures (consistent with the productive-capital channel). The year-on-year change in log NL is therefore the natural outcome variable for Research Question~1.

\textbf{Normalized Difference Built-up Index (NDBI).} Computed from Landsat-8 short-wave infrared and near-infrared bands, NDBI is a continuous index of built-up surface intensity. Whereas NL captures the flow of activity from existing structures, NDBI captures the stock of those structures. A flood that physically damages buildings should reduce NDBI; a flood that only suspends activity without damaging structures should leave NDBI unchanged. This contrast motivates the second outcome equation.

\textbf{Normalized Difference Vegetation Index (NDVI).} NDVI tracks vegetation greenness and is widely used as a proxy for cropland productivity. We use NDVI as a control rather than an outcome, because vegetation responds to floods both negatively (in the immediate inundation period) and positively (in the post-flood season due to silt deposition). Lagged NDVI absorbs these dynamics in the regression specification.

The use of satellite imagery in Indian flood research is itself well established. Jain et al.\ (2006) demonstrated that NOAA AVHRR-derived inundation maps for Bihar and Assam matched ground-truth surveys closely, validating remote sensing as a reliable substitute for fieldwork. The Flood Hazard Atlas of Bihar prepared by NRSC/ISRO and NDMA (2010) operationalised this approach at scale, producing district-level hazard zoning that has since served as the spatial reference for subsequent work. Sinha et al.\ (2008) integrated multiple satellite-derived parameters into a Flood Risk Index for the Kosi basin, establishing that flood vulnerability in Eastern India is multi-dimensional and cannot be captured by a single hydrological variable. The present study draws methodologically on this lineage but extends it in two ways: (i) it uses these proxies to construct outcome variables rather than exposure variables, and (ii) it applies fixed-effects panel estimation rather than cross-sectional zoning.

The flood-exposure variable itself is constructed from the JRC Global Surface Water dataset (Pekel et al., 2016), which provides an internally consistent annual classification of seasonal and permanent water cover at 30-metre resolution.

% =====================================================================
\chapter{Data and Descriptive Statistics}
% =====================================================================

\section{Introduction}
This chapter describes the dataset assembled for the empirical analysis. The unit of observation is the Assembly Constituency (AC) $\times$ year, with the four states of Eastern India --- Bihar, Jharkhand, Odisha, and West Bengal --- observed annually between 2014 and 2019. Section~3.2 sets out the panel structure. Section~3.3 documents the four primary data sources and the construction of each variable. Sections~3.4--3.5 report summary statistics for the regression variables and for flood exposure separately. Section~3.6 presents temporal trends.

\section{Panel Structure}
Table~\ref{tab:panel} summarises the panel structure of the dataset.

\begin{table}[h]
\centering
\caption{Panel structure of the dataset (765 ACs, 108 districts, 2014--2019, $N = 4{,}674$ raw / 3{,}890 estimation).}
\label{tab:panel}
\begin{tabular}{ll}
\toprule
\textbf{Item} & \textbf{Value} \\
\midrule
Unit of observation & Assembly Constituency $\times$ Year \\
States & Bihar, Jharkhand, Odisha, West Bengal \\
Period & 2014--2019 (6 years) \\
Number of ACs & 765 \\
Number of districts (for SE clustering) & 108 \\
Number of AC-year observations (unlagged) & 4{,}674 \\
Estimation sample (with lags; 2014 dropped) & 3{,}890 \\
\bottomrule
\end{tabular}
\end{table}

The panel is a balanced cross-section of 765 Assembly Constituencies observed each year between 2014 and 2019. The 2014 observations are dropped from the estimation sample because the regression specification includes lagged regressors, which are not defined in the first year. AC identifiers are taken from the Election Commission of India delimitation files; district identifiers follow the Census of India 2011 classification.

\section{Data Sources and Variable Construction}
Four primary data sources are combined to construct the panel.

\subsection{Flood Exposure (JRC Global Surface Water)}
The flood-exposure variable, denoted \texttt{Seasonal\_Ratio}, is constructed from the European Commission Joint Research Centre's Global Surface Water dataset (Pekel et al., 2016). For each AC and each year, the dataset is reclassified into four categories --- never water, permanent water, seasonal water, and not-observed --- and the area share of each class is computed via zonal statistics in Google Earth Engine. \texttt{Seasonal\_Ratio} is the share of the AC's land area classified as seasonal water in a given year. It takes values between 0 (no seasonal flooding) and 1 (entire AC under seasonal water), and varies both across ACs and within ACs over time.

\subsection{Nighttime Lights (VIIRS DNB)}
Annual nighttime-light radiance is taken from the NOAA VIIRS Day--Night Band (DNB) composites. For each AC, both the mean and the median radiance across the AC polygon are computed for each year, denoted \texttt{NL\_mean} and \texttt{NL\_median}. The dependent variable in the night-lights regression is the year-on-year change in log radiance, $\Delta \log \text{NL}_{it} = \log(\text{NL}_{it}) - \log(\text{NL}_{i,t-1})$.

\subsection{NDBI and NDVI (Landsat-8)}
The Normalized Difference Built-up Index and the Normalized Difference Vegetation Index are computed from Landsat-8 Surface Reflectance imagery (USGS / Google Earth Engine). NDBI uses the short-wave infrared and near-infrared bands; NDVI uses the near-infrared and red bands. For each AC and each year, both mean and median values across the AC polygon are computed.

\subsection{Auxiliary Identifiers}
Assembly Constituency boundaries and identifiers are taken from the Election Commission of India delimitation files (post-delimitation, valid from 2008). District boundaries follow the Census of India 2011 classification and are used to define the standard-error cluster variable.

\section{Summary Statistics --- Variables}
Table~\ref{tab:sumstats} reports summary statistics at the AC-year level for all regression variables.

\begin{table}[h]
\centering
\caption{Summary statistics for all regression variables, AC-year level.}
\label{tab:sumstats}
\begin{tabular}{lrrrrr}
\toprule
\textbf{Variable} & \textbf{N} & \textbf{Mean} & \textbf{SD} & \textbf{Min} & \textbf{Max} \\
\midrule
Seasonal\_Ratio (flood share)        & 4{,}669 & 0.378  & 0.259  & 0.000   & 0.985 \\
NL --- mean                          & 4{,}674 & 3.019  & 7.432  & 0.054   & 60.258 \\
NL --- median                        & 4{,}674 & 2.632  & 7.363  & 0.044   & 59.562 \\
$\Delta \log$ NL --- mean             & 3{,}895 & 0.126  & 0.497  & $-2.461$ & 4.640 \\
$\Delta \log$ NL --- median           & 3{,}895 & 0.139  & 0.506  & $-2.755$ & 4.891 \\
NDBI --- mean                        & 4{,}674 & $-0.128$ & 0.103 & $-0.447$ & 0.098 \\
NDBI --- median                      & 4{,}674 & $-0.118$ & 0.120 & $-0.490$ & 0.123 \\
NDVI --- mean                        & 4{,}674 & 0.340  & 0.093  & $-0.156$ & 0.664 \\
NDVI --- median                      & 4{,}674 & 0.347  & 0.099  & $-0.343$ & 0.701 \\
\bottomrule
\end{tabular}
\end{table}

Three features of Table~\ref{tab:sumstats} are worth noting. First, the average AC has roughly 38 per cent of its area classified as seasonal water in a given year, with substantial dispersion (SD $= 0.26$) and a maximum approaching unity. Second, the year-on-year change in log NL has a mean of approximately $0.13$ for both aggregations, indicating positive average growth in nighttime-light radiance over the sample period, with substantial variance. Third, NDBI is negative on average (the index is centred around zero by construction; negative values correspond to non-built-up cover), consistent with a predominantly agrarian and rural sample.

\section{Summary Statistics --- Flood Exposure}

\subsection*{By state}

\begin{table}[h]
\centering
\caption{Flood exposure (\texttt{Seasonal\_Ratio}) by state.}
\label{tab:flood-state}
\begin{tabular}{lrrr}
\toprule
\textbf{State} & \textbf{Mean} & \textbf{SD} & \textbf{N} \\
\midrule
Bihar       & 0.323 & 0.198 & 1{,}458 \\
Jharkhand   & 0.547 & 0.311 & \phantom{0}486 \\
Odisha      & 0.444 & 0.265 & \phantom{0}882 \\
West Bengal & 0.341 & 0.238 & 1{,}843 \\
\bottomrule
\end{tabular}
\end{table}

\subsection*{By year}

\begin{table}[h]
\centering
\caption{Flood exposure (\texttt{Seasonal\_Ratio}) by year.}
\label{tab:flood-year}
\begin{tabular}{lrr}
\toprule
\textbf{Year} & \textbf{Mean} & \textbf{SD} \\
\midrule
2014 & 0.034 & 0.050 \\
2015 & 0.412 & 0.217 \\
2016 & 0.470 & 0.234 \\
2017 & 0.472 & 0.223 \\
2018 & 0.407 & 0.227 \\
2019 & 0.474 & 0.221 \\
\bottomrule
\end{tabular}
\end{table}

The state-level breakdown reveals substantial cross-state variation: Jharkhand has the highest average AC-share under seasonal water (0.55), followed by Odisha (0.44), with Bihar and West Bengal lower on average (0.32 and 0.34 respectively). The yearly breakdown shows that 2014 is an outlier with very low flood exposure (mean $= 0.034$). This is a data-availability artefact rather than a true year effect: the JRC dataset's first usable year for our sample begins effectively in 2015. Consequently, 2014 contributes only as a lag year to the estimation sample.

\section{Temporal Trends}
This section presents the time path of each of the eight key variables (mean and median aggregation for each of \texttt{Seasonal\_Ratio}, NL, NDBI, and NDVI) over 2014--2019, averaged across all 765 ACs.

\begin{figure}[h]
\centering
\includegraphics[width=0.78\textwidth]{figures/trend_Seasonal_Ratio_mean.png}
\caption{Average of proportion of land with flooding across Assembly Constituencies, 2014--2019.}
\label{fig:trend-flood-mean}
\end{figure}

\begin{figure}[h]
\centering
\includegraphics[width=0.78\textwidth]{figures/trend_Seasonal_Ratio_median.png}
\caption{Median of proportion of land with flooding across Assembly Constituencies, 2014--2019.}
\label{fig:trend-flood-median}
\end{figure}

\begin{figure}[h]
\centering
\includegraphics[width=0.78\textwidth]{figures/trend_NL_mean.png}
\caption{Average of nighttime lights radiance across Assembly Constituencies, 2014--2019.}
\label{fig:trend-nl-mean}
\end{figure}

\begin{figure}[h]
\centering
\includegraphics[width=0.78\textwidth]{figures/trend_NL_median.png}
\caption{Median of nighttime lights radiance across Assembly Constituencies, 2014--2019.}
\label{fig:trend-nl-median}
\end{figure}

\begin{figure}[h]
\centering
\includegraphics[width=0.78\textwidth]{figures/trend_NDBI_mean.png}
\caption{Average of Normalized Difference Built-up Index across Assembly Constituencies, 2014--2019.}
\label{fig:trend-ndbi-mean}
\end{figure}

\begin{figure}[h]
\centering
\includegraphics[width=0.78\textwidth]{figures/trend_NDBI_median.png}
\caption{Median of Normalized Difference Built-up Index across Assembly Constituencies, 2014--2019.}
\label{fig:trend-ndbi-median}
\end{figure}

\begin{figure}[h]
\centering
\includegraphics[width=0.78\textwidth]{figures/trend_NDVI_mean.png}
\caption{Average of Normalized Difference Vegetation Index across Assembly Constituencies, 2014--2019.}
\label{fig:trend-ndvi-mean}
\end{figure}

\begin{figure}[h]
\centering
\includegraphics[width=0.78\textwidth]{figures/trend_NDVI_median.png}
\caption{Median of Normalized Difference Vegetation Index across Assembly Constituencies, 2014--2019.}
\label{fig:trend-ndvi-median}
\end{figure}

\clearpage
Three patterns emerge from these trends. First, flood exposure (Figures~\ref{fig:trend-flood-mean}--\ref{fig:trend-flood-median}) jumps sharply between 2014 and 2015 as the JRC dataset becomes fully informative, then settles into a stable band of 0.40--0.48 for the rest of the period. Second, NL (Figures~\ref{fig:trend-nl-mean}--\ref{fig:trend-nl-median}) trends gently upward over time, consistent with secular electrification and economic growth. Third, NDBI and NDVI (Figures~\ref{fig:trend-ndbi-mean}--\ref{fig:trend-ndvi-median}) are essentially flat, which is reassuring: it implies that the cross-AC variation we exploit in the regressions is genuine cross-sectional variation rather than a common time trend that would otherwise be absorbed by the year fixed effects.

% =====================================================================
\chapter{Empirical Framework}
% =====================================================================

\section{Introduction}
The theoretical framework in Chapter~2 identifies three channels through which flood exposure affects local economic activity, and Chapter~3 documents the data and variable construction that operationalise those channels. This chapter sets out the two estimating equations used to test the research questions, the variation that identifies each parameter, and the estimation strategy.

\section{Specification --- Equation 1 ($\Delta \log$ NL)}
The first estimating equation models the year-on-year change in log nighttime-light radiance as a function of contemporaneous flood exposure and lagged covariates:
\begin{equation}
\Delta \log \mathrm{NL}_{it} = \beta_1 \, \mathrm{Flood}_{it} + \beta_2 \, \mathrm{NDVI}_{i,t-1} + \beta_3 \, \mathrm{NDBI}_{i,t-1} + \alpha_i + \gamma_t + \varepsilon_{it}.
\label{eq:nl}
\end{equation}
Each term carries a specific interpretation. $\Delta \log \mathrm{NL}_{it}$ is the year-on-year change in log nighttime-light radiance for AC $i$ in year $t$ --- the dependent variable, capturing the flow of economic activity. $\mathrm{Flood}_{it} = \texttt{Seasonal\_Ratio}_{it}$ is contemporaneous flood exposure: the parameter of interest is $\beta_1$. $\mathrm{NDVI}_{i,t-1}$ enters as a lagged control for cropland productivity, and $\mathrm{NDBI}_{i,t-1}$ as a lagged control for built-up intensity. $\alpha_i$ are AC fixed effects, $\gamma_t$ are year fixed effects, and $\varepsilon_{it}$ is the residual.

\section{Specification --- Equation 2 (NDBI)}
The second estimating equation models the level of NDBI as a function of contemporaneous flood exposure, lagged covariates, and the lagged dependent variable:
\begin{equation}
\mathrm{NDBI}_{it} = \beta_1 \, \mathrm{Flood}_{it} + \beta_2 \, \mathrm{NDVI}_{i,t-1} + \beta_3 \, \mathrm{NL}_{i,t-1} + \beta_4 \, \mathrm{NDBI}_{i,t-1} + \alpha_i + \gamma_t + \varepsilon_{it}.
\label{eq:ndbi}
\end{equation}

The structural difference between Equations~(\ref{eq:nl}) and~(\ref{eq:ndbi}) is twofold. First, the dependent variable is the level of NDBI rather than a difference, because the built-up index moves slowly and a year-on-year first-difference would be dominated by measurement noise. Second, the lagged dependent variable $\mathrm{NDBI}_{i,t-1}$ enters on the right-hand side. This captures the fact that built-up intensity at $t$ is heavily determined by built-up intensity at $t-1$.

The inclusion of $\mathrm{NDBI}_{i,t-1}$ together with AC fixed effects produces a small downward bias in the coefficient $\beta_4$, known as Nickell bias. The magnitude of the bias is of order $1/T$, where $T$ is the number of usable time periods. With $T = 5$ in this sample, the bias is approximately 20 per cent. Critically, the flood coefficient $\beta_1$ --- the parameter of interest --- is far less affected by Nickell bias than the lagged-DV coefficient itself, because the flood variable is only weakly correlated with the lagged dependent variable.

\section{Identifying Variation}
The fixed-effects panel estimator exploits two distinct sources of variation in the data.

\textbf{Within-AC variation over time.} The AC fixed effect $\alpha_i$ removes everything that is constant for AC $i$ across the sample period --- terrain, soil type, baseline urbanisation, governance quality, distance to a river. The flood coefficient $\beta_1$ is therefore identified by the variation in flood exposure within an AC across years, not by comparison of one AC to another.

\textbf{Common-shock removal across years.} The year fixed effect $\gamma_t$ removes any shock that affects all ACs in a given year --- monsoon strength, macroeconomic conditions, national policy. Figure~\ref{fig:yearfe} plots the estimated year fixed effects (from a pooled OLS run with AC FE absorbed and 2015 as the omitted year). The pattern is consistent with the macroeconomic record of the period.

\begin{figure}[h]
\centering
\includegraphics[width=0.78\textwidth]{figures/year_fixed_effects_OLS.png}
\caption{Year fixed effects from pooled OLS (AC FE absorbed, year dummies relative to 2015). Common annual shocks of this magnitude would otherwise contaminate the flood coefficient if not absorbed.}
\label{fig:yearfe}
\end{figure}

\textbf{District-clustered standard errors.} Standard errors are clustered at the district level ($n = 108$ districts) to allow for arbitrary correlation between residuals of ACs within the same district --- a natural consequence of weather and flood shocks being district-level events. With 108 clusters, the standard asymptotic inference is comfortably above the Cameron--Miller rule-of-thumb of approximately 50.

\section{Estimation Strategy}
Each of the two equations is estimated under two specifications and two aggregations, yielding four headline regressions per state and a total of twenty regressions across all states.

\textbf{Pooled estimation.} Equations~(\ref{eq:nl}) and~(\ref{eq:ndbi}) are first estimated on the full sample of 3{,}890 AC-year observations across all four states. This produces the average effect of flood exposure across Eastern India.

\textbf{State-by-state estimation.} Each equation is then re-estimated separately for Bihar, Jharkhand, Odisha, and West Bengal. This addresses Research Question~2 directly by allowing $\beta_1$ to differ across states.

\textbf{Mean and median aggregation.} Each model is run twice --- once using the AC-level mean of the satellite raster and once using the median --- as a robustness check on the aggregation choice.

\textbf{Software and reproducibility.} All regressions are estimated in Python using the \texttt{pyfixest} package, which implements two-way fixed-effects estimation with high-dimensional fixed effects and supports cluster-robust inference. Cross-checks against the \texttt{linearmodels} package give numerically identical estimates.

% =====================================================================
\chapter{Results and Discussion}
% =====================================================================

\section{Introduction}
This chapter presents the estimates from the empirical framework set out in Chapter~4. The two estimating equations --- $\Delta \log$ NL on contemporaneous flood and lagged covariates (Equation~\ref{eq:nl}), and NDBI on the same plus the lagged dependent variable (Equation~\ref{eq:ndbi}) --- are each estimated under two aggregations (mean and median) and over five samples (pooled and one per state), yielding twenty regressions in total.

\section{Headline: Forest Plot of Flood Coefficients}
Before proceeding to the individual tables, Figure~\ref{fig:forest} summarises the parameter of interest --- the flood coefficient $\beta_1$ --- across all twenty regressions in a single forest plot.

\begin{figure}[h]
\centering
\includegraphics[width=0.85\textwidth]{figures/forest_plot_flood_OLS.png}
\caption{Estimated effect of flood exposure on $\Delta \log$ NL (left panel) and NDBI (right panel) by state, with 95\% confidence intervals. Pooled estimate at the bottom of each panel.}
\label{fig:forest}
\end{figure}

Three patterns are immediately visible. First, the pooled flood coefficient is statistically indistinguishable from zero in both panels --- a result that would be misleading if reported in isolation. Second, the dispersion across state-level estimates is substantial and the sign varies between states. Third, two estimates stand out as both economically large and statistically significant: the West Bengal NL coefficient (negative, around $-0.32$ to $-0.39$) and the Bihar NDBI coefficient (negative, around $-0.02$).

\section{Night Lights --- Median Model}
Table~\ref{tab:nl-median} reports the median-aggregation estimates of Equation~(\ref{eq:nl}).

\begin{table}[h]
\centering
\caption{Effect of flood exposure on $\Delta \log$ NL (median radiance growth). District-clustered standard errors in parentheses. * $p<0.10$, ** $p<0.05$, *** $p<0.01$.}
\label{tab:nl-median}
\begin{tabular}{lrrrrr}
\toprule
 & \textbf{All states} & \textbf{Bihar} & \textbf{Jharkhand} & \textbf{Odisha} & \textbf{West Bengal} \\
\midrule
Flood (\texttt{Seasonal\_Ratio}) & $-0.061$       & $-0.055$ & $+0.394$ & $+0.235$ & $-0.322$\textsuperscript{**} \\
                                  & (0.083)        & (0.073) & (0.286) & (0.221) & (0.158) \\
NDVI$_{t-1}$                      & $+1.140$\textsuperscript{***} & $+2.556$\textsuperscript{**} & $+8.266$\textsuperscript{**} & $+3.365$ & $+0.804$ \\
                                  & (0.402)        & (1.118) & (3.625) & (2.149) & (0.652) \\
NDBI$_{t-1}$                      & $-1.154$\textsuperscript{**}  & $+1.959$ & $+5.347$ & $-1.113$ & $-3.063$\textsuperscript{***} \\
                                  & (0.557)        & (1.387) & (4.062) & (1.589) & (1.024) \\
\midrule
AC FE      & Yes & Yes & Yes & Yes & Yes \\
Year FE    & Yes & Yes & Yes & Yes & Yes \\
Observations & 3{,}890 & 1{,}215 & 475 & 735 & 1{,}465 \\
R\textsuperscript{2} (within) & 0.027 & 0.041 & 0.082 & 0.038 & 0.052 \\
\bottomrule
\end{tabular}
\end{table}

The pooled flood coefficient ($-0.061$) is small in magnitude and not statistically significant. Three of the four state-level estimates are also insignificant, but West Bengal stands out: the flood coefficient is $-0.322$ ($p<0.05$), indicating that a one-unit increase in the share of an AC under seasonal water is associated with a 32 per cent fall in year-on-year NL growth. This is the largest economically and statistically significant flood effect in the table.

\section{Night Lights --- Mean Model}
Table~\ref{tab:nl-mean} reports the corresponding mean-aggregation estimates.

\begin{table}[h]
\centering
\caption{Effect of flood exposure on $\Delta \log$ NL (mean radiance growth). District-clustered standard errors in parentheses.}
\label{tab:nl-mean}
\begin{tabular}{lrrrrr}
\toprule
 & \textbf{All states} & \textbf{Bihar} & \textbf{Jharkhand} & \textbf{Odisha} & \textbf{West Bengal} \\
\midrule
Flood (\texttt{Seasonal\_Ratio}) & $-0.110$\textsuperscript{*} & $-0.083$\textsuperscript{*} & $+0.439$\textsuperscript{*} & $+0.099$ & $-0.388$\textsuperscript{**} \\
                                  & (0.058)        & (0.046) & (0.231) & (0.197) & (0.166) \\
NDVI$_{t-1}$                      & $+1.298$\textsuperscript{***} & $+3.311$\textsuperscript{***} & $+6.689$\textsuperscript{**} & $+3.949$\textsuperscript{*} & $+0.771$ \\
                                  & (0.472)        & (1.156) & (2.872) & (2.103) & (0.751) \\
NDBI$_{t-1}$                      & $-1.899$\textsuperscript{***} & $+1.830$ & $+3.088$ & $-2.334$ & $-4.054$\textsuperscript{***} \\
                                  & (0.635)        & (1.421) & (3.504) & (1.776) & (1.187) \\
\midrule
AC FE      & Yes & Yes & Yes & Yes & Yes \\
Year FE    & Yes & Yes & Yes & Yes & Yes \\
Observations & 3{,}890 & 1{,}215 & 475 & 735 & 1{,}465 \\
R\textsuperscript{2} (within) & 0.044 & 0.064 & 0.090 & 0.055 & 0.071 \\
\bottomrule
\end{tabular}
\end{table}

The mean specification reinforces the median results and sharpens them slightly. West Bengal's flood coefficient is now $-0.388$ ($p<0.05$), Bihar's flood coefficient becomes marginally significant at the 10 per cent level ($\beta = -0.083$, $p<0.10$), and Jharkhand's flood coefficient flips to a marginally positive estimate ($+0.439$, $p<0.10$) which is best interpreted as a small-sample artefact (Jharkhand has only 475 AC-year observations).

\section{NDBI --- Median Model}
Table~\ref{tab:ndbi-median} reports the median-aggregation estimates of Equation~(\ref{eq:ndbi}).

\begin{table}[h]
\centering
\caption{Effect of flood exposure on NDBI (median, level). The lagged-DV coefficient (NDBI$_{t-1}$) is subject to Nickell bias of order $1/T \approx 20\%$; the flood coefficient is largely unaffected.}
\label{tab:ndbi-median}
\begin{tabular}{lrrrrr}
\toprule
 & \textbf{All states} & \textbf{Bihar} & \textbf{Jharkhand} & \textbf{Odisha} & \textbf{West Bengal} \\
\midrule
Flood (\texttt{Seasonal\_Ratio}) & $-0.016$        & $-0.023$\textsuperscript{***} & $-0.052$ & $+0.019$\textsuperscript{*} & $+0.010$ \\
                                  & (0.010)        & (0.007) & (0.033) & (0.011) & (0.016) \\
NDVI$_{t-1}$                      & $+0.002$        & $+0.346$\textsuperscript{***} & $-0.208$\textsuperscript{***} & $+0.127$\textsuperscript{**} & $-0.015$ \\
                                  & (0.031)        & (0.066) & (0.071) & (0.061) & (0.033) \\
NL$_{t-1}$                        & $+0.0013$\textsuperscript{***} & $+0.0059$\textsuperscript{***} & $+0.0007$ & $+0.0031$\textsuperscript{***} & $+0.0007$\textsuperscript{*} \\
                                  & (0.0004)       & (0.0009) & (0.0005) & (0.0010) & (0.0004) \\
NDBI$_{t-1}$                      & $-0.058$\textsuperscript{*}  & $+0.108$\textsuperscript{*} & $+0.026$ & $-0.029$ & $+0.015$ \\
                                  & (0.030)        & (0.060) & (0.051) & (0.040) & (0.044) \\
\midrule
AC FE      & Yes & Yes & Yes & Yes & Yes \\
Year FE    & Yes & Yes & Yes & Yes & Yes \\
Observations & 3{,}890 & 1{,}215 & 475 & 735 & 1{,}465 \\
R\textsuperscript{2} (within) & 0.020 & 0.199 & 0.227 & 0.043 & 0.011 \\
\bottomrule
\end{tabular}
\end{table}

The pooled flood effect on NDBI is small and insignificant, but Bihar shows a highly significant negative effect ($\beta = -0.023$, $p<0.01$). This is consistent with the productive-capital channel set out in Chapter~2: floods in Bihar physically damage the built-up stock, lowering the NDBI level even after controlling for the lagged dependent variable. Odisha shows a small positive effect ($\beta = +0.019$, $p<0.10$), which may reflect post-flood reconstruction activity rather than damage. The lagged NL coefficient is positive and significant in four of the five samples, indicating that higher prior economic activity is associated with subsequent built-up expansion --- a sensible cross-validation of the two satellite series.

\section{NDBI --- Mean Model}
Table~\ref{tab:ndbi-mean} reports the corresponding mean-aggregation estimates.

\begin{table}[h]
\centering
\caption{Effect of flood exposure on NDBI (mean, level). District-clustered standard errors in parentheses.}
\label{tab:ndbi-mean}
\begin{tabular}{lrrrrr}
\toprule
 & \textbf{All states} & \textbf{Bihar} & \textbf{Jharkhand} & \textbf{Odisha} & \textbf{West Bengal} \\
\midrule
Flood (\texttt{Seasonal\_Ratio}) & $-0.013$        & $-0.021$\textsuperscript{***} & $-0.042$ & $+0.015$ & $+0.010$ \\
                                  & (0.009)        & (0.006) & (0.027) & (0.009) & (0.013) \\
NDVI$_{t-1}$                      & $+0.004$        & $+0.282$\textsuperscript{***} & $-0.142$\textsuperscript{**} & $+0.104$\textsuperscript{**} & $-0.011$ \\
                                  & (0.026)        & (0.062) & (0.052) & (0.044) & (0.027) \\
NL$_{t-1}$                        & $+0.0012$\textsuperscript{***} & $+0.0050$\textsuperscript{***} & $+0.0007$\textsuperscript{*} & $+0.0022$\textsuperscript{***} & $+0.0007$\textsuperscript{**} \\
                                  & (0.0003)       & (0.0008) & (0.0004) & (0.0007) & (0.0003) \\
NDBI$_{t-1}$                      & $-0.065$\textsuperscript{**} & $+0.051$ & $+0.086$\textsuperscript{**} & $-0.047$ & $+0.019$ \\
                                  & (0.032)        & (0.060) & (0.038) & (0.030) & (0.048) \\
\midrule
AC FE      & Yes & Yes & Yes & Yes & Yes \\
Year FE    & Yes & Yes & Yes & Yes & Yes \\
Observations & 3{,}890 & 1{,}215 & 475 & 735 & 1{,}465 \\
R\textsuperscript{2} (within) & 0.024 & 0.195 & 0.212 & 0.047 & 0.017 \\
\bottomrule
\end{tabular}
\end{table}

The mean specification confirms the pattern from Table~\ref{tab:ndbi-median}. Bihar's flood coefficient remains negative and significant at the 1 per cent level ($\beta = -0.021$), with magnitude almost identical to the median estimate. The Odisha effect loses significance once the mean is used, but its sign remains positive. The lagged NL coefficient is significant across all four state regressions, reinforcing its interpretation as a robust predictor of built-up expansion.

\section{Robustness: Mean vs Median Aggregation}
A natural concern is whether the choice of aggregating the satellite raster by mean versus median affects the conclusions. Figure~\ref{fig:meanvsmedian} plots the flood coefficient under both aggregations side-by-side.

\begin{figure}[h]
\centering
\includegraphics[width=0.85\textwidth]{figures/mean_vs_median_flood_OLS.png}
\caption{Robustness of the flood coefficient under mean versus median aggregation. Panel A: $\Delta \log$ NL outcome. Panel B: NDBI outcome. Error bars are 95\% confidence intervals. Significance stars: * $p<0.10$, ** $p<0.05$, *** $p<0.01$.}
\label{fig:meanvsmedian}
\end{figure}

Two observations follow. First, the sign of the flood coefficient is identical under both aggregations in every state and for both outcomes --- a meaningful robustness check. Second, the magnitude is generally larger under mean aggregation than under median aggregation, because the mean is more sensitive to bright outliers (urban centres, illuminated highways) that the flood disrupts disproportionately. The qualitative findings --- significant negative effect in West Bengal for NL, significant negative effect in Bihar for NDBI, no effect elsewhere --- are not driven by the aggregation choice.

\section{Summary: Heatmap of All Twenty Regressions}
To facilitate cross-comparison, Figure~\ref{fig:heatmap} presents the full set of coefficient estimates from all twenty regressions in a single heatmap.

\begin{figure}[h]
\centering
\includegraphics[width=0.95\textwidth]{figures/heatmap_all_coefficients_OLS.png}
\caption{Coefficient summary of all twenty OLS regressions. Rows are state $\times$ aggregation; columns are regressors. Cells display the estimated coefficient with significance stars. The colour scale encodes coefficient sign (red = negative, blue = positive).}
\label{fig:heatmap}
\end{figure}

Three secondary patterns visible in Figure~\ref{fig:heatmap} are worth noting. First, NDVI$_{t-1}$ is consistently positive and significant in the NL regressions for Bihar, Jharkhand, and Odisha, indicating that better pre-existing vegetation (a proxy for cropland productivity) is associated with stronger subsequent NL growth --- consistent with the agricultural-land channel of Chapter~2. Second, NDVI$_{t-1}$ flips sign in Jharkhand's NDBI regressions (negative and significant), which likely reflects the fact that Jharkhand's most vegetated areas are forested rather than agricultural, and forested cover does not translate into built-up expansion. Third, NL$_{t-1}$ loads positively and significantly across all four state-level NDBI regressions, providing internal cross-validation between the two satellite series.

\section{Discussion of Cross-State Heterogeneity}
The most striking feature of the results is the cross-state heterogeneity revealed by state-by-state estimation, which would be entirely missed if only the pooled regressions were reported.

\textbf{West Bengal --- economic-activity channel.} The flood effect operates almost exclusively through nighttime lights, with a magnitude of 32 to 39 per cent of a one-unit increase in seasonal-water share. NDBI is unaffected. This is consistent with West Bengal's economic structure: a densely populated, partially urbanised state with extensive electrical infrastructure where floods disrupt activity (the disruption channel) without permanently destroying built-up stock.

\textbf{Bihar --- built-up infrastructure channel.} The flood effect operates through NDBI, with a small but precisely estimated decline of approximately 2 per cent of an NDBI unit per unit of seasonal-water share. NL is at most marginally affected. This is consistent with Bihar's pattern of recurrent breach-induced flooding from Himalayan-fed rivers, which physically damages housing and roads (the productive-capital channel) more than it disrupts electrical activity.

\textbf{Jharkhand and Odisha --- no consistent effect.} Neither outcome shows a consistently significant flood effect in these two states. For Jharkhand, the small sample size (475 AC-years, only 81 ACs) limits statistical power; the marginally positive NL coefficient under mean aggregation should be treated as an artefact rather than a meaningful estimate. For Odisha, the absence of effect may reflect a genuinely smaller flood-damage profile under the period of analysis, or the offsetting role of post-flood reconstruction activity (suggested by the marginally positive NDBI estimate).

\textbf{Implication for pooled estimation.} The four-state pooled regressions yield insignificant flood coefficients precisely because the state-level effects partially cancel out in the pooled sample: West Bengal's strong negative NL effect is offset by Jharkhand's positive estimate, and Bihar's NDBI effect is diluted by the absence of an effect elsewhere. The pooled coefficient is therefore not the right summary statistic for policy purposes. State-by-state estimation is essential for both diagnosis and design.

% =====================================================================
\chapter{Summary and Concluding Remarks}
% =====================================================================

\section{Summary of Findings}
This study set out to estimate the short-run economic impact of flood exposure on Assembly Constituencies in Eastern India, using a panel of 765 ACs across Bihar, Jharkhand, Odisha, and West Bengal observed annually between 2014 and 2019. Two outcome variables were modelled --- the year-on-year change in log nighttime-light radiance (a proxy for the flow of economic activity) and the level of the Normalized Difference Built-up Index (a proxy for the stock of built-up infrastructure) --- using OLS with Assembly Constituency and year fixed effects, and standard errors clustered at the district level.

Four headline findings emerge:
\begin{enumerate}[label=(\roman*),leftmargin=2.2em,itemsep=4pt]
\item \textbf{Pooling masks the effect.} When the four states are pooled, the flood coefficient is statistically indistinguishable from zero for both outcomes. This is not because floods have no economic impact in Eastern India, but because the state-level effects partially offset each other in the pooled sample.
\item \textbf{West Bengal: a significant economic-activity contraction.} Flood exposure significantly reduces nighttime-lights growth in West Bengal, with a coefficient of approximately $-0.32$ (median, $p<0.05$) to $-0.39$ (mean, $p<0.05$). NDBI is largely unaffected, suggesting the impact operates through the disruption channel rather than the productive-capital channel.
\item \textbf{Bihar: a significant built-up-infrastructure decline.} Flood exposure significantly reduces NDBI in Bihar ($\beta \approx -0.02$, $p<0.01$). The nighttime-lights effect is at most marginally significant. This pattern is consistent with the productive-capital channel --- recurrent breach-induced flooding from Himalayan-fed rivers physically damages structures rather than only disrupting activity.
\item \textbf{Jharkhand and Odisha: no consistent effect.} Neither state shows a robustly significant flood coefficient on either outcome, which may reflect smaller sample sizes, topographic differences, or genuine heterogeneity in damage exposure.
\end{enumerate}

The robustness check using mean versus median aggregation confirms that these findings are not artefacts of the aggregation choice; the sign of every flood coefficient is preserved across both aggregations.

\section{Policy Implications}
The state-by-state heterogeneity identified above carries direct policy implications.

A uniform flood-mitigation policy is unlikely to be efficient. Flood-management budgets in Eastern India are currently allocated through state-level engineering interventions --- embankments, barrages, and drainage works --- that implicitly assume a common damage profile. The estimates in Chapter~5 suggest that this assumption is empirically incorrect: the channels through which floods damage local economies differ by state, and so should the policy response.

\textbf{West Bengal} requires interventions oriented toward protecting economic activity during floods --- back-up power supply for commercial centres, faster road clearance, mobile market infrastructure --- because the binding constraint there is short-run disruption rather than long-run capital damage.

\textbf{Bihar} requires interventions oriented toward protecting and repairing built-up infrastructure --- flood-resilient housing standards, faster post-flood reconstruction, and embankment maintenance --- because the binding constraint is the physical destruction of standing structures.

\textbf{Jharkhand and Odisha} require further study before specific policy prescriptions can be made, but the absence of a measurable effect under the present specification does not imply that floods are economically harmless in these states. It only implies that the channels through which they operate are not adequately captured by the proxies used here.

The broader implication is that flood-mitigation budgets should be designed against state-level damage profiles, not a national average.

\section{Caveats and Limitations}
Several caveats apply to the interpretation of the results.

\textbf{Nickell bias.} The NDBI specification includes the lagged dependent variable on the right-hand side together with AC fixed effects. This produces a small downward bias on the lagged-DV coefficient of order $1/T$. With $T = 5$ usable years, the bias is approximately 20 per cent. The flood coefficient --- the parameter of interest --- is far less affected, but the magnitude of the lagged NDBI coefficient should be interpreted with caution.

\textbf{Missing road-connectivity control.} Road density is plausibly correlated with both flood exposure (low-lying areas may have less developed road networks) and economic activity. The Geo Sadak road-network data is not currently incorporated into the specification; including it as a control would strengthen identification, and is the most important direction for refining the present analysis.

\textbf{Short panel.} The estimation sample contains five usable years (2015--2019). This limits the precision of state-level estimates, particularly for Jharkhand (475 AC-years), and constrains the use of more demanding identification strategies such as long-difference designs or anticipation-effect tests.

\textbf{Satellite proxies versus ground truth.} Nighttime lights, NDBI, and NDVI are useful precisely because they are observable at high frequency and fine spatial resolution, but they remain proxies. Each captures only a partial dimension of the local economy, and the mapping from satellite signal to economic outcome is necessarily indirect. Validation against household survey data --- where it becomes available at AC level --- would strengthen confidence in the estimates.

\section{Future Work}
Three extensions follow naturally from the present analysis.

\textbf{(i) Flood $\times$ road-connectivity interaction.} Using AC-level road-density data from the Geo Sadak road-network database, the regression can be augmented with an interaction term between flood exposure and road density. The hypothesis is that better-connected ACs experience smaller economic disruption during floods because evacuation, supply chains, and post-flood recovery all operate faster on well-connected road networks.

\textbf{(ii) Urban--rural heterogeneity.} The damage profile of a flood is likely to differ between predominantly urban and predominantly rural ACs. By interacting flood exposure with the baseline NDBI level (or by splitting the sample at the median NDBI), the present specification can be extended to test whether the same flood event has differentially larger effects on built-up versus non-built-up ACs.

\textbf{(iii) Channel decomposition with disaggregated outcomes.} The current specification estimates two reduced-form effects (on activity and on built-up stock). A natural extension is to decompose these into the underlying channels --- for instance, by combining the present panel with data on road damage, agricultural yields, and labour mobility, allowing the agricultural-land, productive-capital, and disruption channels to be quantified separately.

These extensions can be implemented within the same panel-fixed-effects framework developed in this thesis and represent a coherent research agenda for further empirical work on the economic impacts of flooding in Eastern India.

% =====================================================================
% REFERENCES (chapter*, listed in TOC)
% =====================================================================
\clearpage
\addcontentsline{toc}{chapter}{References}
\chapter*{References}
\markboth{References}{References}

\noindent\textbf{Academic papers}
\begin{description}[leftmargin=2.5em,style=nextline,itemsep=2pt]
\item[--] Jain, S.~K., Saraf, A.~K., Goswami, A.\ \& Ahmad, T.\ (2006). Flood inundation mapping using NOAA AVHRR data. \textit{Natural Hazards}, 37, 43--59.
\item[--] Kale, V.~S.\ (2003). Geomorphic effects of monsoon floods on Indian rivers. \textit{Natural Hazards}, 28, 1--3.
\item[--] Nickell, S.\ (1981). Biases in dynamic models with fixed effects. \textit{Econometrica}, 49(6), 1417--1426.
\item[--] Pekel, J.-F., Cottam, A., Gorelick, N.\ \& Belward, A.~S.\ (2016). High-resolution mapping of global surface water and its long-term changes. \textit{Nature}, 540, 418--422.
\item[--] Sinha, R., Bapalu, G.~V., Singh, L.~K.\ \& Rath, B.\ (2008). Flood risk analysis in the Kosi river basin, North Bihar using multi-parametric approach of analytical hierarchy process (AHP). \textit{Journal of the Indian Society of Remote Sensing}, 36(2), 293--307.
\end{description}

\noindent\textbf{Reports and government publications}
\begin{description}[leftmargin=2.5em,style=nextline,itemsep=2pt]
\item[--] Census of India (2011). \textit{Population and Housing Census of India}. Office of the Registrar General and Census Commissioner, Government of India.
\item[--] Central Water Commission (CWC) (2024). \textit{Annual Flood Damage Report}. Ministry of Jal Shakti, Government of India.
\item[--] Election Commission of India (2008). \textit{Delimitation Files for Assembly Constituencies (post-delimitation)}. Election Commission of India, New Delhi.
\item[--] Government of Bihar (2024). \textit{State Flood Reports}. Department of Disaster Management, Government of Bihar.
\item[--] Ministry of Statistics and Programme Implementation (MoSPI) (2024). \textit{State-wise Net State Domestic Product per capita}. Government of India.
\item[--] National Disaster Management Authority (NDMA) (2024). \textit{Flood Frequency and Damage Statistics for India}. Government of India.
\item[--] National Flood Commission (1980). \textit{Report on Flood-Prone Area Statistics for India}. Government of India.
\item[--] National Institute of Disaster Management (NIDM) (2024). \textit{Flood Hazard Atlas of India}. Ministry of Home Affairs, Government of India.
\item[--] National Remote Sensing Centre (NRSC), ISRO, NDMA \& Government of Bihar (2010). \textit{Flood Hazard Atlas of Bihar --- A Geospatial Approach}. NRSC-ISRO, Hyderabad.
\end{description}

\noindent\textbf{Datasets}
\begin{description}[leftmargin=2.5em,style=nextline,itemsep=2pt]
\item[--] Google Earth Engine (2024). \textit{Landsat-8 Surface Reflectance imagery}. United States Geological Survey (USGS) / NASA.
\item[--] Joint Research Centre, European Commission (2024). \textit{Global Surface Water Dataset}. European Commission Joint Research Centre, Ispra. (Built on Pekel et al., 2016.)
\item[--] Ministry of Road Transport and Highways (MoRTH) (2024). \textit{Geo Sadak --- Road Infrastructure Database}. Government of India.
\item[--] NOAA / Visible Infrared Imaging Radiometer Suite (VIIRS) (2024). \textit{Day-Night Band (DNB) Annual Composites}. National Oceanic and Atmospheric Administration, USA.
\item[--] NOAA / Defense Meteorological Satellite Program (DMSP) (2014). \textit{Operational Linescan System Nighttime Lights}. National Oceanic and Atmospheric Administration, USA.
\end{description}

\noindent\textbf{Software}
\begin{description}[leftmargin=2.5em,style=nextline,itemsep=2pt]
\item[--] Berge, L., Krantz, S.\ \& McDermott, G.\ (2024). \textit{pyfixest: Fast high-dimensional fixed-effects regression in Python}. Open-source Python package.
\item[--] Sheppard, K.\ (2024). \textit{linearmodels: Linear models for Python (panel and instrumental-variables estimation)}. Open-source Python package.
\end{description}

% =====================================================================
% APPENDIX
% =====================================================================
\appendix
\renewcommand{\thechapter}{\Alph{chapter}}
\titleformat{\chapter}[display]
  {\normalfont\Large\bfseries\centering}
  {APPENDIX \thechapter}
  {12pt}
  {\Large\MakeUppercase}

\chapter{Variable Construction Details}

\section{Flood Exposure (\texttt{Seasonal\_Ratio})}
The flood-exposure variable is built from the Joint Research Centre's Global Surface Water dataset (Pekel et al., 2016), accessed via Google Earth Engine. For each calendar year, the dataset provides a per-pixel classification of water occurrence at 30-metre resolution. Pixels are categorised into four classes: 0 = never water, 1 = seasonal water, 2 = permanent water, and 3 = not observed.

For each Assembly Constituency (AC) polygon, the area-weighted share of each class is computed as:
\[
\texttt{Seasonal\_Ratio}_{it} = \frac{\text{HISTO\_1}_{it}}{\text{HISTO\_0}_{it} + \text{HISTO\_1}_{it} + \text{HISTO\_2}_{it}}.
\]
The denominator excludes the not-observed class to avoid attributing missing data to either land or water. \texttt{Seasonal\_Ratio} takes values in $[0, 1]$ and is the primary measure of contemporaneous flood exposure used throughout the empirical chapters.

\section{Nighttime Lights (VIIRS DNB)}
Annual nighttime-light radiance is taken from the NOAA VIIRS Day--Night Band annual composites (\texttt{NOAA/VIIRS/DNB/MONTHLY\_V1}), aggregated to annual means. For each AC polygon and each year, two summary statistics are computed via zonal statistics: \texttt{NL\_mean}$_{it}$ (pixel-wise mean radiance within AC $i$, year $t$) and \texttt{NL\_median}$_{it}$ (pixel-wise median).

The dependent variable in Equation~(\ref{eq:nl}) is the year-on-year change in log radiance, computed within each AC: $\Delta \log \text{NL}_{it} = \log(\text{NL}_{it}) - \log(\text{NL}_{i,t-1})$. Pre-2014 DMSP-OLS data are not used because they are not directly comparable with VIIRS without intercalibration; for that reason, 2014 is the start year of the panel.

\section{NDBI and NDVI (Landsat-8 SR)}
The Normalized Difference Built-up Index and the Normalized Difference Vegetation Index are computed at the pixel level from Landsat-8 Surface Reflectance imagery (\texttt{LANDSAT/LC08/C02/T1\_L2}):
\[
\text{NDBI} = \frac{\text{SWIR1} - \text{NIR}}{\text{SWIR1} + \text{NIR}}, \qquad
\text{NDVI} = \frac{\text{NIR} - \text{RED}}{\text{NIR} + \text{RED}},
\]
where SWIR1 is band 6 (short-wave infrared), NIR is band 5 (near-infrared), and RED is band 4. Annual cloud-free composites are constructed by median-compositing all valid pixels within each calendar year.

\section{Zonal Aggregation}
All raster-to-polygon aggregations are performed using zonal-statistics functions in Google Earth Engine (\texttt{ee.Image.reduceRegions}). The polygon source is the Election Commission of India post-delimitation Assembly Constituency boundary file (effective from 2008). Reducer parameters are set to compute mean, median, and pixel count for every (AC, year) combination, ensuring identical spatial aggregation across the four satellite indices used in the analysis.

\section{Lag Construction}
Lagged regressors are constructed within each Assembly Constituency by sorting on year and shifting by one period. For any variable $X$:
\[
X_{i,t-1} = X_{i, t-1} \quad \text{if both } (i, t-1) \text{ and } (i, t) \text{ are observed; missing otherwise}.
\]
This procedure drops the 2014 observations from the estimation sample because no lag is defined for the panel's first year. The estimation sample therefore contains 5 usable years (2015--2019).

\section{Cluster Variable for Standard Errors}
Standard errors in all regressions are clustered at the district level. The cluster identifier is constructed by concatenating the state code (\texttt{ST\_CODE}) and the district code (\texttt{DT\_CODE}) from the Census of India 2011 classification, yielding a unique two-digit-pair identifier per district. The estimation sample contains 108 unique district clusters --- comfortably above the Cameron--Miller rule of thumb of approximately 50 clusters needed for reliable cluster-robust inference.

\end{document}
